\subsection{Exception Handling Infrastructure}

Python exception handling is built around a structured trap. The \texttt{try} block marks a region where exceptional control flow may begin. The \texttt{except} blocks describe which exception classes can be intercepted. The optional \texttt{else} block describes code that should run only when the protected region succeeds. The optional \texttt{finally} block describes code that must run whether the protected region succeeds, fails, or transfers control outward.

The full structure is:

\begin{verbatim}
try:
    protected_statements
except SomeError:
    recovery_statements
else:
    success_only_statements
finally:
    always_run_statements
\end{verbatim}

Not every part is required in every program. The most common form is \texttt{try}/\texttt{except}. The other clauses exist because exception handling is not only about catching failures. It is also about separating normal success logic from recovery logic and ensuring cleanup logic runs unconditionally.

\begin{quote}
\texttt{try} marks the danger region. \texttt{except} handles selected failures. \texttt{else} runs only after success. \texttt{finally} runs no matter how the region is left.
\end{quote}

\subsubsection{The \texttt{try} / \texttt{except} Trap: Intercepting and Deflecting Specific Exception Class Trees}

A \texttt{try}/\texttt{except} statement creates a controlled region where specific exceptional routes can be intercepted. If an exception is raised inside the \texttt{try} block, Python stops the ordinary route through that block and searches the attached \texttt{except} clauses for a matching exception class.

\begin{verbatim}
text = "forty-two"

try:
    number = int(text)
    print(f"number: {number}")
except ValueError:
    print("conversion failed")

print("after try/except")
\end{verbatim}

Possible output:

\begin{verbatim}
conversion failed
after try/except
\end{verbatim}

The line:

\begin{verbatim}
print(f"number: {number}")
\end{verbatim}

does not run, because \texttt{int(text)} raises \texttt{ValueError}. Execution is deflected into the matching \texttt{except ValueError:} block.

The control route is:

\begin{verbatim}
enter try
    -> attempt int(text)
        success:
            bind number
            print number
            continue after try/except
        ValueError:
            jump to except ValueError
            print conversion failed
            continue after try/except
\end{verbatim}

The \texttt{try} block should normally contain only the statements that belong to the risky operation and its immediate success-dependent continuation. If too much unrelated code is placed inside the same \texttt{try} block, the handler may catch exceptions from a larger region than intended.

Too broad:

\begin{verbatim}
text = "42"

try:
    number = int(text)
    print(f"number: {number}")
    missing_name = unknown_variable
except ValueError:
    print("conversion failed")
\end{verbatim}

Possible output:

\begin{verbatim}
number: 42
Traceback (most recent call last):
  ...
NameError: name 'unknown_variable' is not defined
\end{verbatim}

The \texttt{NameError} is not caught because the handler catches only \texttt{ValueError}. That is good: the handler should not pretend that a programmer mistake is an input conversion problem.

If the handler were too broad, it could hide the real issue:

\begin{verbatim}
text = "42"

try:
    number = int(text)
    print(f"number: {number}")
    missing_name = unknown_variable
except Exception as err:
    print("some ordinary exception was caught")
    print(f"type(err): {type(err)}")
    print(f"message: {err}")
\end{verbatim}

Possible output:

\begin{verbatim}
number: 42
some ordinary exception was caught
type(err): <class 'NameError'>
message: name 'unknown_variable' is not defined
\end{verbatim}

The program did not have a conversion problem. It had a name-resolution problem. Catching \texttt{Exception} trapped both possibilities.

A better structure keeps the risky conversion narrow:

\begin{verbatim}
text = "42"

try:
    number = int(text)
except ValueError as err:
    print("conversion failed")
    print(f"message: {err}")
else:
    print(f"number: {number}")
    print("conversion succeeded")
\end{verbatim}

Possible output:

\begin{verbatim}
number: 42
conversion succeeded
\end{verbatim}

The handler now corresponds exactly to the operation that can raise the expected \texttt{ValueError}.

Multiple handlers can be attached to one \texttt{try} block:

\begin{verbatim}
items = ["10", "20", "D'oh!"]
index = 2

try:
    text = items[index]
    number = int(text)
    print(f"number: {number}")
except IndexError as err:
    print("bad index")
    print(f"type(err): {type(err)}")
except ValueError as err:
    print("bad integer text")
    print(f"type(err): {type(err)}")
    print(f"message: {err}")
\end{verbatim}

Possible output:

\begin{verbatim}
bad integer text
type(err): <class 'ValueError'>
message: invalid literal for int() with base 10: "D'oh!"
\end{verbatim}

The first risky operation is indexing. It may raise \texttt{IndexError}. The second risky operation is conversion. It may raise \texttt{ValueError}. The handlers separate the two exceptional routes.

The exception object is an ordinary runtime object:

\begin{verbatim}
try:
    number = int("No soup for you!")
except ValueError as err:
    print(f"type(err): {type(err)}")
    print(f"dir(err) sample: {dir(err)[:12]}")
    print(f"err.args: {err.args}")
    print(f"str(err): {str(err)}")
    print(f"is Exception: {isinstance(err, Exception)}")
\end{verbatim}

Possible output:

\begin{verbatim}
type(err): <class 'ValueError'>
dir(err) sample: ['__cause__', '__class__', '__context__', '__delattr__',
'__dict__', '__dir__', '__doc__', '__eq__', '__format__', '__ge__',
'__getattribute__', '__getstate__']
err.args: ('invalid literal for int() with base 10: 'No soup for you!'',)
str(err): invalid literal for int() with base 10: 'No soup for you!'
is Exception: True
\end{verbatim}

The \texttt{args} attribute stores the arguments carried by the exception object. The \texttt{str()} conversion gives the message. The \texttt{isinstance()} test confirms that the object belongs to the exception class tree.

The \texttt{try}/\texttt{except} structure is therefore a class-based routing trap:

\begin{quote}
If the protected route raises an exception whose class matches a handler, Python deflects execution into that handler instead of continuing through the ordinary route.
\end{quote}

\subsubsection{The Exception \texttt{else} Clause: Isolating Code Paths That Run Only When No Exceptions Occur}

The \texttt{else} clause of a \texttt{try} statement runs only if the \texttt{try} block finishes without raising an exception. It is not the same as an \texttt{if}-\texttt{else} branch. It is a success-only continuation attached to exception handling.

\begin{verbatim}
text = "42"

try:
    number = int(text)
except ValueError as err:
    print("conversion failed")
    print(f"message: {err}")
else:
    print("conversion succeeded")
    print(f"number: {number}")
\end{verbatim}

Possible output:

\begin{verbatim}
conversion succeeded
number: 42
\end{verbatim}

If the conversion fails, the \texttt{else} block is skipped:

\begin{verbatim}
text = "forty-two"

try:
    number = int(text)
except ValueError as err:
    print("conversion failed")
    print(f"message: {err}")
else:
    print("conversion succeeded")
    print(f"number: {number}")
\end{verbatim}

Possible output:

\begin{verbatim}
conversion failed
message: invalid literal for int() with base 10: 'forty-two'
\end{verbatim}

The route is:

\begin{verbatim}
try block succeeds:
    -> skip except
    -> run else

try block raises handled exception:
    -> run matching except
    -> skip else
\end{verbatim}

The main reason to use \texttt{else} is to keep the protected region narrow. Code in the \texttt{else} block depends on success, but it is not part of the risky operation being guarded.

Less precise:

\begin{verbatim}
text = "42"

try:
    number = int(text)
    print(f"number: {number}")
    print("Nobody expects the Spanish Inquisition!")
except ValueError:
    print("conversion failed")
\end{verbatim}

More precise:

\begin{verbatim}
text = "42"

try:
    number = int(text)
except ValueError:
    print("conversion failed")
else:
    print(f"number: {number}")
    print("Nobody expects the Spanish Inquisition!")
\end{verbatim}

Possible output:

\begin{verbatim}
number: 42
Nobody expects the Spanish Inquisition!
\end{verbatim}

The second version says exactly what is being protected: the conversion. The printing after successful conversion is success-path logic, not conversion-risk logic.

The same pattern is useful with dictionary lookup:

\begin{verbatim}
scores = {
    "Jerry": 72,
    "Elaine": 91,
    "Kramer": 42,
}

name = "Kramer"

try:
    score = scores[name]
except KeyError as err:
    print(f"missing score for: {err}")
else:
    print(f"{name:<8} | {score:>3}")
    print("score lookup succeeded")
\end{verbatim}

Possible output:

\begin{verbatim}
Kramer   |  42
score lookup succeeded
\end{verbatim}

If the key is absent:

\begin{verbatim}
scores = {
    "Jerry": 72,
    "Elaine": 91,
    "Kramer": 42,
}

name = "Newman"

try:
    score = scores[name]
except KeyError as err:
    print(f"missing score for: {err}")
else:
    print(f"{name:<8} | {score:>3}")
    print("score lookup succeeded")
\end{verbatim}

Possible output:

\begin{verbatim}
missing score for: 'Newman'
\end{verbatim}

The dictionary object has its own lookup behavior:

\begin{verbatim}
scores = {
    "Jerry": 72,
    "Elaine": 91,
    "Kramer": 42,
}

print(f"type(scores): {type(scores)}")
print(f"dir(scores) sample: {dir(scores)[:12]}")
print(f"has __getitem__: {'__getitem__' in dir(scores)}")
print(f"has get: {'get' in dir(scores)}")
print(f"has __contains__: {'__contains__' in dir(scores)}")
\end{verbatim}

Possible output:

\begin{verbatim}
type(scores): <class 'dict'>
dir(scores) sample: ['__class__', '__class_getitem__', '__contains__',
'__delattr__', '__delitem__', '__dir__', '__doc__', '__eq__',
'__format__', '__ge__', '__getattribute__', '__getitem__']
has __getitem__: True
has get: True
has __contains__: True
\end{verbatim}

The bracket operation \texttt{scores[name]} uses lookup behavior associated with \texttt{\_\_getitem\_\_()}. A missing key raises \texttt{KeyError}. The \texttt{else} block runs only if that lookup succeeds.

The practical rule is:

\begin{quote}
Put only the risky operation in \texttt{try}. Put the code that should run after successful completion in \texttt{else}.
\end{quote}

This keeps exception routes readable. The \texttt{except} block describes failure recovery. The \texttt{else} block describes the ordinary success continuation.

\subsubsection{Cleansing and Post-Processing: The Unconditional Execution Invariants of the \texttt{finally} Block}

The \texttt{finally} block runs after the \texttt{try} structure is left, regardless of whether the protected code succeeded, failed, or was interrupted by a control-flow transfer such as \texttt{break}, \texttt{return}, or an uncaught exception.

At this stage of the chapter, the simplest rule is:

\begin{quote}
\texttt{finally} means: run this cleanup code no matter how the protected region was exited.
\end{quote}

Success case:

\begin{verbatim}
text = "42"

try:
    number = int(text)
    print(f"number: {number}")
except ValueError:
    print("conversion failed")
finally:
    print("finally block ran")

print("after try statement")
\end{verbatim}

Possible output:

\begin{verbatim}
number: 42
finally block ran
after try statement
\end{verbatim}

Handled failure case:

\begin{verbatim}
text = "forty-two"

try:
    number = int(text)
    print(f"number: {number}")
except ValueError:
    print("conversion failed")
finally:
    print("finally block ran")

print("after try statement")
\end{verbatim}

Possible output:

\begin{verbatim}
conversion failed
finally block ran
after try statement
\end{verbatim}

In both programs, the \texttt{finally} block runs.

The \texttt{finally} block also runs when the exception is not handled locally. To demonstrate that without stopping the entire script, an outer \texttt{try} can catch the propagated exception:

\begin{verbatim}
try:
    try:
        print("inner try begins")
        number = int("D'oh!")
        print(f"number: {number}")
    finally:
        print("inner finally ran before propagation")
except ValueError as err:
    print("outer handler caught the propagated exception")
    print(f"type(err): {type(err)}")
    print(f"message: {err}")
\end{verbatim}

Possible output:

\begin{verbatim}
inner try begins
inner finally ran before propagation
outer handler caught the propagated exception
type(err): <class 'ValueError'>
message: invalid literal for int() with base 10: "D'oh!"
\end{verbatim}

The inner \texttt{try} has no \texttt{except} handler. The \texttt{ValueError} must propagate outward. But before that propagation reaches the outer handler, the inner \texttt{finally} block runs.

This means \texttt{finally} is not a success block and not a failure block. It is an exit block.

A typical use is cleanup state:

\begin{verbatim}
resource_open = False

try:
    print("opening resource")
    resource_open = True

    print("using resource")
    number = int("42")
    print(f"number: {number}")
finally:
    print("cleanup section")
    resource_open = False
    print(f"resource_open: {resource_open}")
\end{verbatim}

Possible output:

\begin{verbatim}
opening resource
using resource
number: 42
cleanup section
resource_open: False
\end{verbatim}

If an error happens, cleanup still runs:

\begin{verbatim}
resource_open = False

try:
    print("opening resource")
    resource_open = True

    print("using resource")
    number = int("forty-two")
    print(f"number: {number}")
except ValueError as err:
    print("conversion failed")
    print(f"message: {err}")
finally:
    print("cleanup section")
    resource_open = False
    print(f"resource_open: {resource_open}")
\end{verbatim}

Possible output:

\begin{verbatim}
opening resource
using resource
conversion failed
message: invalid literal for int() with base 10: 'forty-two'
cleanup section
resource_open: False
\end{verbatim}

The variable \texttt{resource\_open} is only a simple teaching flag here. In real programs, the cleanup section commonly closes files, releases locks, closes network connections, ends temporary states, or restores invariants.

A \texttt{finally} block can also run when a loop is interrupted with \texttt{break}:

\begin{verbatim}
items = ["soup", "pretzels", "muffin"]

for item in items:
    try:
        print(f"checking: {item}")

        if item == "pretzels":
            print("pretzels found")
            break
    finally:
        print(f"finished attempt for: {item}")

print("after loop")
\end{verbatim}

Possible output:

\begin{verbatim}
checking: soup
finished attempt for: soup
checking: pretzels
pretzels found
finished attempt for: pretzels
after loop
\end{verbatim}

The \texttt{break} leaves the loop, but the \texttt{finally} block for the current iteration still runs before execution reaches the statement after the loop.

The full clause order is fixed:

\begin{verbatim}
try:
    ...
except SomeError:
    ...
else:
    ...
finally:
    ...
\end{verbatim}

The \texttt{else} clause, when present, comes after all \texttt{except} clauses and before \texttt{finally}. The \texttt{finally} block is last because it runs after the success or failure route has been determined.

All three routes together:

\begin{verbatim}
texts = ["42", "forty-two"]

for text in texts:
    print(f"input: {text!r}")

    try:
        number = int(text)
    except ValueError as err:
        print("except route")
        print(f"message: {err}")
    else:
        print("else route")
        print(f"number: {number}")
    finally:
        print("finally route")
        print("-" * 20)
\end{verbatim}

Possible output:

\begin{verbatim}
input: '42'
else route
number: 42
finally route
--------------------
input: 'forty-two'
except route
message: invalid literal for int() with base 10: 'forty-two'
finally route
--------------------
\end{verbatim}

For the first input, the \texttt{try} block succeeds, so \texttt{else} runs, then \texttt{finally}. For the second input, the \texttt{try} block raises \texttt{ValueError}, so \texttt{except} runs, then \texttt{finally}.

The practical reading is:

\begin{itemize}
    \item \texttt{try}: attempt the protected operation;
    \item \texttt{except}: recover from selected exception classes;
    \item \texttt{else}: run success-only code when no exception occurred;
    \item \texttt{finally}: run cleanup code regardless of success or failure.
\end{itemize}

For control-flow graph analysis, \texttt{finally} adds a guaranteed exit checkpoint. Ordinary success routes, handled exception routes, unhandled propagation routes, and local loop-interruption routes must pass through it before continuing outward.